\documentclass[11pt]{article}
\usepackage[utf8]{inputenc}
\usepackage{booktabs}
\usepackage{graphicx}
\usepackage{amsmath}
\usepackage{hyperref}
\usepackage{geometry}
\geometry{margin=1in}

\title{Metabolic State Analysis Using User-Supplied Metabolic Genes}
\author{QA\_R5 Research Task}
\date{\today}

\begin{document}

\maketitle

\section{Overview}

This report presents the analysis of metabolic gene expression patterns using the user-supplied gene set \texttt{metabolic\_genes\_total.csv} applied to the Choudhury et al. 2022 meningioma single-cell RNA-seq dataset (3CA:20773). The analysis addresses three questions: (1) whether metabolic genes alone reveal reproducible transcriptional metabolic groupings across all cells, (2) the relationship between these groupings and known cell types/subtypes, and (3) whether distinct metabolic groupings exist within the \texttt{cell\_subtype=CD8 T cells} population.

\section{Data and Methods}

\subsection{Dataset}

The analysis uses the Choudhury et al. 2022 meningioma dataset from the Curated Cancer Cell Atlas (3CA:20773) \cite{choudhury_2022}. The dataset contains 58,843 cells and 33,538 genes from 10 meningioma samples. User-supplied metabolic genes were provided in \texttt{inputs/metabolic\_genes\_total.csv} with 1,988 source symbols.

\subsection{Gene Matching}

The gene matching process identified 1,646 matched genes from the 1,988 user-supplied symbols. Of these, 341 matched via exact symbol match and 1,457 via case-insensitive match. 190 symbols remained unmatched. Duplicate expression symbols were summed in raw counts before QC/normalization.

\subsection{Normalization and Preprocessing}

Cells were filtered for QC passing criteria. Expression was normalized using per-cell library size scaling to 10,000 UMI equivalents followed by log1p transformation:

\begin{equation}
    x_{\text{norm}} = \log_1\left( \frac{\sum_{g} x_{g,c}}{N_{\text{cells}}} \times 10,000 + 1 \right)
\end{equation}

where $x_{g,c}$ is the raw UMI count for gene $g$ in cell $c$, and $N_{\text{cells}}$ is the total number of cells. Features were scaled using Scanpy's \texttt{scale} function with \texttt{zero\_center=False} and \texttt{max\_value=10}.

\subsection{Clustering}

Clustering was performed using Scanpy's Leiden algorithm with resolution 0.8 for the global analysis and resolution 0.4 for the CD8 T cells subset. The resolution was selected based on highest median silhouette across seeds, with mean pairwise ARI as tie-breaker. Stability was assessed via Leiden randomization on a fixed PCA/neighbor graph.

\section{Question 1: Reproducible Metabolic Groupings Across All Cells}

\subsection{Results}

The global analysis of 58,843 cells using 1,646 metabolic genes yielded 17 clusters at resolution 0.8. Cluster sizes ranged from 714 to 8,806 cells.

\begin{itemize}
    \item \textbf{Leiden silhouette:} 0.217
    \item \textbf{Matched-k KMeans:} 0.220
    \item \textbf{Random control comparison:} 2 of 19 random controls exceeded metabolic clustering (add-one rank fraction: 0.15)
    \item \textbf{HVG vs random controls:} HVG silhouette (0.283) exceeded random control range [0.199, 0.222]
\end{itemize}

\subsection{Interpretation}

The metabolic gene clustering achieved a silhouette of 0.217, comparable to matched-k KMeans (0.220) but lower than HVG-based clustering (0.283). The random control comparison shows that 2 of 19 random gene sets exceeded the metabolic clustering in silhouette (add-one rank fraction: 0.15), indicating the metabolic clustering is not significantly better than random. The HVG control demonstrates that using highest-variance genes yields stronger clustering (0.283) than the metabolic gene set.

\section{Question 2: Relationship to Cell Types and Subtypes}

\subsection{Results}

\begin{table}[h!]
\centering
\caption{Confounder NMI values for global analysis}
\label{tab:confounders}
\begin{tabular}{lrr}
\toprule
Confounder & Global NMI & Within-replicate range \\
\midrule
Patient & 0.518 & 0.071 -- 0.676 \\
Sample & 0.548 & -- \\
Cell type & 0.496 & 0.379 -- 0.745 \\
Cell subtype & 0.648 & 0.224 -- 0.680 \\
\bottomrule
\end{tabular}
\end{table}

The categorical confounder NMI analysis reveals substantial association between clustering and cell type/subtype. Cell subtype shows the strongest association (NMI: 0.648), followed by patient (NMI: 0.518) and sample (NMI: 0.548). Cell type shows moderate association (NMI: 0.496).

\subsection{Interpretation}

The clustering is strongly associated with cell subtype (NMI: 0.648) and patient (NMI: 0.518), suggesting that the metabolic groupings are largely driven by cell type identity and sample/patient effects rather than true metabolic states. The within-replicate NMI ranges show substantial variation, with the weakest replicate (patient 3) having NMI as low as 0.224. This indicates that the clustering is not robustly reproducible within biological replicates.

\section{Question 3: Metabolic Groupings Within CD8 T Cells}

\subsection{Results}

The CD8 T cells subset analysis included 3,624 cells and 1,318 metabolic genes, yielding 6 clusters at resolution 0.4. Cluster sizes ranged from 55 to 2,002 cells.

\begin{itemize}
    \item \textbf{Leiden silhouette:} 0.099
    \item \textbf{Matched-k KMeans:} 0.095
    \item \textbf{Random control comparison:} 17 of 19 random controls exceeded metabolic clustering (add-one rank fraction: 0.9)
    \item \textbf{HVG vs random controls:} HVG silhouette (0.168) exceeded random control range [0.092, 0.131]
\end{itemize}

\subsection{Interpretation}

The CD8 T cells analysis shows very weak clustering structure (silhouette: 0.099), substantially lower than the global analysis. The random control comparison is particularly concerning: 17 of 19 random controls exceeded the metabolic clustering (add-one rank fraction: 0.9), indicating the metabolic clustering is worse than random. The HVG control shows that HVG-based clustering (0.168) is stronger than the metabolic gene set.

\section{Conclusions}

\subsection{Question 1 Answer}

The metabolic genes do not form strongly reproducible transcriptional metabolic groupings across all cells. The clustering silhouette (0.217) is comparable to matched-k KMeans (0.220) but lower than HVG-based clustering (0.283). The random control comparison (2/19 controls exceeded, add-one rank: 0.15) indicates the metabolic clustering is not significantly better than random.

\subsection{Question 2 Answer}

The metabolic groupings are strongly associated with cell type and subtype (cell subtype NMI: 0.648, cell type NMI: 0.496), as well as patient (NMI: 0.518) and sample (NMI: 0.548). This suggests the groupings are driven by cell identity and sample effects rather than true metabolic states. The within-replicate NMI ranges show substantial variation, indicating limited reproducibility within biological replicates.

\subsection{Question 3 Answer}

Within CD8 T cells, the metabolic gene clustering is very weak (silhouette: 0.099) and worse than random (17/19 controls exceeded, add-one rank: 0.9). The HVG control shows HVG-based clustering (0.168) is stronger than the metabolic gene set. There is no evidence of distinct metabolic groupings within CD8 T cells.

\section{Limitations}

\begin{itemize}
    \item Cell-level clustering metrics are descriptive and do not establish discrete biological states.
    \item The analysis treats cells as nested within patients/samples; cell-level NMI, ARI, silhouette, and cluster composition are descriptive statistics.
    \item Random control comparisons use add-one rank fraction as a descriptive metric, not a biological P value.
    \item The clustering resolution was selected based on silhouette and ARI criteria, not biological validation.
    \item No automatic threshold proves independence from confounders.
    \item The same matrix-derived cluster labels cannot be used for cell-level significance testing.
\end{itemize}

\begin{thebibliography}{9}

\bibitem{choudhury_2022}
Abrar Choudhury, Stephen T. Magill, Charlotte D. Eaton, et al.
\textit{Meningioma DNA methylation groups identify biological drivers and therapeutic vulnerabilities}.
\textbf{Nature Genetics}, 54(5):649-659, 2022.
\url{https://doi.org/10.1038/s41588-022-01061-8}

\bibitem{gavish_2023}
Avishai Gavish, Michael Tyler, Alissa C. Greenwald, et al.
\textit{Hallmarks of transcriptional intratumour heterogeneity across a thousand tumours}.
\textbf{Nature}, 618(7965):598-606, 2023.
\url{https://doi.org/10.1038/s41586-023-06130-4}

\end{thebibliography}

\end{document}
